\section{Ordered logarithms and positive modular calculus}
\label{sec:positive-calculus}

The logarithm is where commutative and noncommutative expansions part company.
For a scalar gas it selects connected clusters. For an operator gas it must
also remember the order of mutually incompatible factors. Heaps retain that
order economically. Their coefficients explain both the cancellations that
make a logarithm local in a formal sense and the obstruction to estimating it
by absolute scalar cluster weights. A different representation, based on
positive random times, then gives an operator approximation with a useful
error certificate.

\subsection{Heaps, order polynomials, and the logarithm}

Fix a finite alphabet with a symmetric reflexive dependence relation. Two
distinct letters commute precisely when they are independent. A word modulo
these commutations is a \emph{heap}: its occurrences form a partially ordered
set, with the order generated by dependent letters appearing in word order.
This is the Cartier--Foata--Viennot description of partial commutation
\citep{CartierFoata1969,Viennot1986}. It records the information that would be
lost by replacing ordered products by commutative monomials.

Let $M=\sum_C(-1)^{|C|}C$ and $Z=\sum_CC$, where $C$ runs over cliques of
mutually independent letters, including the empty clique. Both are formal
series in the completed heap algebra. For a nonempty heap $H$, write
$\Omega_H(t)$ for the polynomial counting weakly order-preserving maps from
$H$ to $\{1,\ldots,t\}$.

\begin{proposition}[The order-polynomial logarithm]
\label{op:heap}
For every nonempty heap,
\begin{equation}
 [H](-\log M)=\Omega_H'(0),\qquad
 [H]\log Z=(-1)^{|H|+1}\Omega_H'(0).
 \label{op:heap-coefficient}
\end{equation}
Disconnected heaps have zero coefficient. For a natural labeling of a heap
with $n$ elements, the first coefficient also equals
\begin{equation}
 \Omega_H'(0)=\sum_{\pi\in\mathcal L(H)}
 \frac{(-1)^{\operatorname{des}(\pi)}}
 {n\binom{n-1}{\operatorname{des}(\pi)}},
 \label{op:eulerian}
\end{equation}
where $\mathcal L(H)$ denotes its linear extensions.
\end{proposition}

\begin{proof}
For a positive integer $t$, the coefficient of $H$ in $M^t$ is
$(-1)^n\overline\Omega_H(t)$: choosing the clique in each factor is a strict
order-preserving map. Polynomial reciprocity gives
$\overline\Omega_H(t)=(-1)^n\Omega_H(-t)$. Differentiating at $t=0$ yields
the first identity; substituting the negative of each letter gives the
second. A disjoint union has a product of order polynomials, each vanishing
at zero, so its derivative vanishes. Finally, the standard linear-extension
expansion
\[
 \Omega_H(t)=\sum_{\pi\in\mathcal L(H)}
 \binom{t+n-1-\operatorname{des}(\pi)}n
\]
gives \eqref{op:eulerian} by differentiation.
\end{proof}

There are two useful interpretations of the coefficient. If
$\mathcal O(H)\subset[0,1]^H$ is Stanley's order polytope, then
$\Omega_H(t)=L_{\mathcal O(H)}(t-1)$, where $L$ is its Ehrhart polynomial
\citep{Stanley1986}. Thus the logarithm evaluates a shifted lattice-point
polynomial at a negative argument. Formula \eqref{op:eulerian}, on the other
hand, is the Eulerian-idempotent expression familiar in Lie-theoretic
logarithms and their pre-Lie refinements \citep{BandieraSchaetz2017}. These
are structural identifications of the coefficient, rather than positivity
statements about it.

For example, let $H_m$ consist of $m$ independent leaves below one root.
Then $\Omega_{H_m}(t)=\sum_{j=1}^tj^m$ and
\begin{equation}
 \Omega_{H_m}'(0)=B_m(1),
 \label{op:bernoulli}
\end{equation}
with the Bernoulli-polynomial convention $B_1(1)=1/2$. The coefficients
change sign and vanish for odd $m>1$. A connected heap need not make a
positive contribution to an ordered logarithm.

\fig{fig07_ordered_heaps}{Ordered logarithms retain signed combinatorics.
(a) A star heap consists of independent leaves below a common root; its order
polynomial is a sum of powers. (b) The derivatives at zero are Bernoulli
values; positive and negative nonzero values are distinguished by color,
and gray crosses indicate zeros. The vertical axis records the logarithm of
the magnitude only for nonzero coefficients. (c) The root-linear logarithmic
response of the positive operator gas in the text has Taylor radius
$2\sin(\pi/m)$ for $m\ge7$. It eventually lies below the fixed leaf activity
$1/16$, despite a uniform scalar cluster criterion.}{fig:heaps}

\paragraph{A positive gas for which absolute ordered expansion fails.}
On $m$ qubits put $P_i=|1\rangle\langle1|_i$ and
$X=|0^m\rangle\langle1^m|+|1^m\rangle\langle0^m|$. Consider
\[
 Z_m=\prod_{i=1}^m(I+a_iP_i)+bX,
 \qquad a_i=a=\frac1{16},\qquad b=\frac1{8\,4^m}.
\]
Its least eigenvalue is positive: the diagonal product is at least $I$ and
$\|bX\|=b<1$. Treating each singleton and the full support as polymers,
the Koteck\'y--Preiss condition with size cost $|S|\log2$ remains valid even
after the additional support tilt $2^{|S|}$. Indeed, the incompatible sum at
a leaf is $4a+4^mb=3/8\le\log2$, and at the root is
$4ma+4^mb=m/4+1/8\le m\log2$.

Nevertheless the $(0^m,1^m)$ matrix element of the derivative of $\log Z_m$
at $b=0$, with all leaf activities equal, is
\begin{equation}
 f_m(a)=\frac{m\log(1+a)}{(1+a)^m-1}.
 \label{op:root-response}
\end{equation}
This is the divided difference of the scalar logarithm between the diagonal
eigenvalues $1$ and $(1+a)^m$. For $m\ge7$ its nearest nonremovable poles
are at $a=e^{\pm2\pi i/m}-1$, with modulus $2\sin(\pi/m)<1$.
At $m=128$ this is smaller than $1/16$. Thus even the root-linear part of
the absolute multivariate Taylor expansion cannot converge at the displayed
positive activities: absolute convergence would imply convergence on the
diagonal. The scalar criterion controls its own commutative expansion; the
ordered logarithm requires a cancellation-aware representation.

\subsection{The normalization of a quantum response}
\label{op:normalization}

Let $\rho>0$ be a density matrix, $K=\log\rho$, and
$\dot\rho=\dot\rho^*$ a trace-zero tangent. Two natural Hermitian inputs are
\[
 A=\rho^{-1/2}\dot\rho\rho^{-1/2},\qquad
 H=\frac12(\dot\rho\rho^{-1}+\rho^{-1}\dot\rho).
\]
They are related but not interchangeable. Write $D=\ad_K$. Divided
differences in an eigenbasis of $\rho$ give
\begin{equation}
 D\log_\rho(\dot\rho)
 =\frac{D/2}{\sinh(D/2)}A
 =\frac D{\sinh D}H.
 \label{op:log-response}
\end{equation}
For instance, the first multiplier at eigenvalues $p_i,p_j$ is
$\sqrt{p_ip_j}(\log p_i-\log p_j)/(p_i-p_j)$, with diagonal value one.
The half-argument in the first expression is forced by the relative
normalization. It is relevant when translating between modular response,
quantum scores, and numerical implementations.

The functions in \eqref{op:log-response} are characteristic functions.
With the Fourier convention $f(x)=\int e^{itx}p(t)\,dt$,
\begin{equation}
 \frac{x/2}{\sinh(x/2)}
 =\int_{\R}e^{itx}\frac{\pi}{2\cosh^2(\pi t)}\,dt,
 \qquad
 \frac x{\sinh x}
 =\int_{\R}e^{itx}\frac{\pi}{4\cosh^2(\pi t/2)}\,dt.
 \label{op:fourier-log}
\end{equation}
Consequently each multiplier is a positive average of unitary
conjugations $A\mapsto e^{itK}Ae^{-itK}$. This replaces ordered Taylor
coefficients by an exact completely positive integral. The scalar infinite
divisibility has a classical probabilistic theory \citep{PitmanYor2003};
positive-definite modular multipliers also belong to the theory of
quantum information metrics \citep{HiaiKosakiPetzRuskai2013}.

\subsection{Gamma clocks and errors measured on the observable}

For a Hermitian $G$, let $L=\ad_G^2$. The heat map $e^{-sL}$ is the average
of $A\mapsto e^{itG}Ae^{-itG}$ for a centered Gaussian time of variance
$2s$. It is unital, trace preserving, completely positive, and an operator
norm contraction. We use this representation even though $L$ is a
superoperator rather than a positive map.

\begin{theorem}[Positive hyperbolic clock and local error]
\label{op:hyperbolic-clock}
For $\tau\ge0$,
\begin{equation}
 T_\tau=\left(\frac{\ad_G}{\sinh\ad_G}\right)^\tau
 =\mathbb E e^{-S_\tau L},\qquad
 S_\tau=\sum_{j\ge1}Y_j,\quad
 Y_j\sim\operatorname{Gamma}\left(\tau,\text{rate }\pi^2j^2\right)
 \label{op:gamma-clock}
\end{equation}
with independent summands. It is a random-unitary semigroup. Put
$S_{\tau,R}=\sum_{j\le R}Y_j$ and
$a_R=\mathbb E(S_\tau-S_{\tau,R})$. The finite and mean-corrected channels
obey
\begin{align}
 \|\mathbb Ee^{-S_{\tau,R}L}A-T_\tau A\|
 &\le\frac{\tau}{\pi^2R}\|\ad_G^2A\|,\label{op:clock-first}\\
 \|\mathbb Ee^{-(S_{\tau,R}+a_R)L}A-T_\tau A\|
 &\le\frac{\tau}{6\pi^4R^3}\|\ad_G^4A\|.
 \label{op:clock-mean}
\end{align}
\end{theorem}

\begin{proof}
Euler's product for $\sinh$ and the Laplace transform of a Gamma variable
give the identity. The sum exists almost surely because its mean is
$\tau\sum_j(\pi j)^{-2}=\tau/6$. For the first bound, integrate the
generator along the omitted nonnegative time and use contraction. For the
second, write the omitted time as $U$. Taylor's theorem with integral
remainder around $\mathbb EU$ is valid on every segment between these
nonnegative times. Its linear term has mean zero, and its remainder is at
most $\tfrac12\operatorname{Var}(U)\|L^2A\|$. The independent Gamma
variances give
\[
 \operatorname{Var}(U)=\tau\sum_{j>R}(\pi j)^{-4}
 \le\frac{\tau}{3\pi^4R^3}.
\]
The retained heat channel is a contraction, which completes both estimates.
\end{proof}

These bounds involve the supplied observable and its actual commutators.
Replacing them by $\|G\|$ would discard useful spatial information and
could introduce the total system size. Conversely, the relative input $A$
in \eqref{op:log-response} need not have the same support as $\dot\rho$.
The theorem does not assert spatial locality of modular evolution.

\fig{fig08_modular_clocks}{Positive modular filters and finite random-time
approximations. (a) The normalized power-response multiplier
$\sinh(rx)/(r\sinh x)$ approaches $x/\sinh x$ as $r\downarrow0$.
(b) Its positive Fourier density is
$\sin(\pi r)/[2r(\cosh(\pi t)+\cos(\pi r))]$; each curve integrates to one.
(c) At scalar frequency $4$ and semigroup time $1$, retaining $R$ Euler
factors has an $R^{-1}$ error bound, while replacing the omitted clock by
its mean gives an $R^{-3}$ bound. Solid curves are exact scalar errors and
dashed curves are the stated analytic bounds; this panel tests the scalar
multiplier, not a many-body interaction norm.}{fig:clocks}

\subsection{Bessel functions connect the clock to quantum information geometry}
\label{op:bessel}

The hyperbolic example belongs to a larger family. For $\nu>-1$, define
\[
 \mathcal I_\nu(x)=\Gamma(\nu+1)(2/x)^\nu I_\nu(x)
 ={}_0F_1(\nu+1;x^2/4),\qquad \phi_\nu(x)=\mathcal I_\nu(x)^{-1}.
\]
The value at zero is one. If $j_{\nu,k}$ are the positive zeros of $J_\nu$,
then the Bessel product gives
\[
 \phi_\nu(x)=\prod_{k\ge1}(1+\lambda_kx^2)^{-1},\qquad
 \lambda_k=j_{\nu,k}^{-2},\qquad
 \sum_k\lambda_k=\frac1{4(\nu+1)}.
\]
Thus $\phi_\nu(\ad_G)^\tau$ is a random-unitary semigroup with clock
$\sum_k\operatorname{Gamma}(\tau,\text{scale }\lambda_k)$. At $\tau=1$
this is one half of the exit time from radius one of a Bessel process of
dimension $2\nu+2$, started at zero. The reciprocal-product and probability
connections are classical; finite moment-optimal Gamma approximations are
developed by Clark \citep{Clark2026}. Here their operator use is governed by
the commutator estimates below.

To connect with a quantum statistical model, set $G=K/2$ and
\begin{equation}
 m_\nu(t)=\sqrt t\,\mathcal I_\nu(\tfrac12\log t).
 \label{op:bessel-mean}
\end{equation}
A symmetric normalized operator-monotone function $m$ defines the Petz
information metric through
\[
 \mathcal J_{m,\rho}=R_\rho m(L_\rho R_\rho^{-1}),\qquad
 g_{m,\rho}(\dot\rho_i,\dot\rho_j)
 =\Re\tr\bigl(\dot\rho_i\mathcal J_{m,\rho}^{-1}(\dot\rho_j)\bigr),
\]
where $L_\rho$ and $R_\rho$ mean left and right multiplication
\citep{Petz1996,HiaiKosakiPetzRuskai2013}. Operator monotonicity ensures
that information cannot increase after a quantum channel.

\begin{proposition}[Exact metric range]
\label{op:metric-range}
The function \eqref{op:bessel-mean} defines a symmetric normalized
operator-monotone mean exactly for $\nu\ge-1/2$. In that range,
\begin{equation}
 \mathcal J_{m_\nu,\rho}^{-1}(\dot\rho)
 =\phi_\nu(\ad_K/2)A.
 \label{op:bessel-score}
\end{equation}
The endpoint is the symmetric logarithmic derivative (SLD) metric;
$\nu=1/2$ is the Bogoliubov--Kubo--Mori (BKM) metric; the limit as
$\nu\to\infty$ is the geometric-mean metric.
\end{proposition}

\begin{proof}
For $\nu>-1/2$ the Bessel integral representation becomes
\[
 m_\nu(t)=\mathbb E t^U,\qquad
 U\sim\operatorname{Beta}(\nu+\tfrac12,\nu+\tfrac12).
\]
It is a positive mixture of the operator-monotone functions $t^u$,
$0\le u\le1$; the symmetric beta law gives $tm_\nu(1/t)=m_\nu(t)$.
The endpoint law puts mass $1/2$ at each of $0,1$. Accordingly,
\[
 m_{-1/2}(t)=\frac{1+t}2,\qquad
 m_{1/2}(t)=\frac{t-1}{\log t},\qquad
 m_\infty(t)=\sqrt t.
\]
For $-1<\nu<-1/2$ direct differentiation gives
$m_\nu''(1)=-(2\nu+1)/[8(\nu+1)]>0$, contradicting the scalar concavity
of an operator-monotone function. Finally, comparing entries in an
eigenbasis of $\rho$ proves \eqref{op:bessel-score}.
\end{proof}

This has an operational distinction worth preserving. SLD Fisher
information describes local parameter estimation; the BKM metric is the
affine Hessian of $\tr(\rho\log\rho)$ and the quadratic term of quantum
relative entropy. The remaining means interpolate between information
geometries. Complete positivity of their normalized score filters holds
even below $\nu=-1/2$, but monotonicity of the corresponding metric does
not. For example, $\nu=-3/4$, $\rho=\diag(4/5,1/5)$ and
$\dot\rho=X/10$ give information approximately $0.033246$, whereas
measurement in the $X$ basis gives classical Fisher information $0.04$.
This numerical illustration is supplementary; the exact second-derivative
test already proves the boundary of the metric family.

\subsection{Positive quadrature produces information-matrix brackets}
\label{op:quadrature}

An infinite clock is not needed for computation. Let
$\mu=\sum_k\lambda_k\delta_{\lambda_k}$, a finite positive measure on a
bounded interval, and write
\[
 g(L)=\int F_\lambda(L)\,d\mu(\lambda),\qquad
 F_\lambda(L)=\frac{\log(I+\lambda L)}\lambda
 =\int_0^1 L(I+u\lambda L)^{-1}\,du.
\]
At $\lambda=0$ the continuous value is $L$. Approximating this generator
by positive quadrature retains the entire channel semigroup. It also makes
two complementary error directions available.

\begin{theorem}[Finite positive quadrature with local certificates]
\label{op:quadrature-theorem}
Let $\mu$ be a finite positive measure of bounded support in $[0,\infty)$,
with the nondegenerate moment matrices needed for the following rules.
For its $r$-node Gaussian rule $(c_i,w_i)$, put
\[
 T_r^G(\tau)=\prod_{i=1}^r(I+c_iL)^{-\tau w_i/c_i},\qquad
 h_r=\int p_r(\lambda)^2\,d\mu(\lambda),
\]
where $p_r$ is the monic orthogonal polynomial. Then
\begin{equation}
 \|(T_r^G(\tau)-e^{-\tau g(L)})A\|
 \le\frac{\tau h_r}{2r+1}\|\ad_G^{4r+2}A\|.
 \label{op:gauss-error}
\end{equation}
For the Gaussian rule $(c_i,v_i)$ of $\lambda\,d\mu$, set
$w_i=v_i/c_i$ and $w_0=\mu([0,\infty))-\sum_iw_i\ge0$. The Radau rule
\[
 T_r^R(\tau)=e^{-\tau w_0L}
 \prod_{i=1}^r(I+c_iL)^{-\tau w_i/c_i}
\]
satisfies, with $\widetilde p_r$ monic orthogonal for $\lambda\,d\mu$,
\begin{equation}
 \|(T_r^R(\tau)-e^{-\tau g(L)})A\|
 \le\frac{\tau\widetilde h_r}{2r+2}\|\ad_G^{4r+4}A\|,
 \quad
 \widetilde h_r=\int\lambda\widetilde p_r(\lambda)^2\,d\mu.
 \label{op:radau-error}
\end{equation}
All maps displayed are random-unitary channels for $\tau\ge0$. Their
scalar multipliers obey $T_r^R(\tau;s)\le e^{-\tau g(s)}\le T_r^G(\tau;s)$
for $s\ge0$.
\end{theorem}

\begin{proof}
Differentiating the resolvent representation gives
\[
 \partial_\lambda^kF_\lambda(L)
 =(-1)^kk!\int_0^1u^k L^{k+1}(I+u\lambda L)^{-k-1}\,du.
\]
Resolvent powers are positive averages of heat channels and commute with
$L$. Thus their action on $A$ has norm at most
$k!\|L^{k+1}A\|/(k+1)$. The Gaussian Hermite-interpolation remainder has
factor $p_r^2$ and derivative order $2r$; the Radau remainder has factor
$\lambda\widetilde p_r^2$ and derivative order $2r+1$.
Integration therefore bounds the generator errors by the right sides of
\eqref{op:gauss-error}--\eqref{op:radau-error} without $\tau$.
For the operator-valued remainder this statement follows from the positive
Peano-kernel integral for Hermite interpolation, followed by the norm bound;
equivalently, scalarize by norm-one functionals and apply Hahn--Banach.
The signed scalar derivative formula gives the stated multiplier directions.
Commuting-semigroup Duhamel and contraction multiply the generator bounds
by $\tau$. Gaussian and Radau weights are positive, so each approximant
is a finite Gamma clock, with an additional deterministic heat time for
Radau. Degenerate finite-support cases terminate exactly.
An atom at zero contributes its continuous heat factor
$\exp(-\tau w_0L)$ rather than a Gamma factor with zero scale.
\end{proof}

For Bessel clocks the Rayleigh sums satisfy
\begin{equation}
 s_1=\frac1{4(\nu+1)},\qquad
 (\nu+n)s_n=\sum_{i=1}^{n-1}s_i s_{n-i},\quad n\ge2.
 \label{op:rayleigh}
\end{equation}
This follows by inserting the logarithmic derivative of the Bessel product
into its differential equation. Rational $\nu$ gives exact rational moments
and hence algebraically specified quadrature nodes and weights. For the
hyperbolic clock $\nu=1/2$, one Gaussian node already gives
\begin{equation}
 T_1^G(\tau)=(I+L/15)^{-5\tau/2},\qquad
 \|\mathrm{error}(A)\|\le\frac{\tau}{9450}\|\ad_G^6A\|.
 \label{op:one-clock}
\end{equation}
The supplemental implementation isolates the quadrature nodes by rational
Sturm intervals and encloses the weights independently of the floating
eigensolver. A separate rounding budget is added before a numerical rule
is called certified.

For a statistical model, let
$\langle A,B\rangle_{\rho,1/2}=\Re\tr(\rho^{1/2}A^*\rho^{1/2}B)$.
Since $\ad_K^2$ is self-adjoint and nonnegative for this inner product,
the scalar multiplier inequalities become finite information-matrix bounds:
\begin{equation}
 \mathsf G_r^R\preceq\mathsf G_\nu\preceq\mathsf G_r^G,
 \qquad
 (\mathsf G_r^*)_{ij}=
 \langle A_i,T_r^*(1)A_j\rangle_{\rho,1/2}.
 \label{op:information-bracket}
\end{equation}
This is Loewner order on matrices of parameter information. It is not
positive-map order: the difference of distinct unital maps cannot be
positive, since a positive map vanishing at the identity vanishes everywhere.
Nor is each finite approximation necessarily a monotone metric for all
states. Equation \eqref{op:information-bracket} is the fixed-state promise
needed for certified statistical calculations.

\fig{fig12_information_clocks}{Bessel clocks connect positive channels to
quantum information metrics. (a) The symmetric beta laws yield the exact
operator means $m_\nu(t)=\mathbb E[t^U]$; the displayed densities have
$\nu=0,1/2,2$. (b) The relative-score filters include the symmetric
logarithmic derivative (SLD), the Bogoliubov--Kubo--Mori (BKM) score, and
the limit associated with the geometric mean. (c) One positive Gaussian
quadrature node gives the upper scalar multiplier and one Radau node plus
heat drift gives the lower multiplier, enclosing $x/\sinh x$. The shaded
interval gives fixed-state information-matrix bounds through functional
calculus; it does not express an order between completely positive maps.
All curves are evaluated from the stated formulas.}{fig:information-clocks}

\paragraph{A positive empirical output.}
Here $G=K/2$, so modular unitary conjugation is an isometry for
$\langle\cdot,\cdot\rangle_{\rho,1/2}$.
Let $B_i=T_r^*(1/2)A_i$. The desired information matrix is the Gram matrix
of the $B_i$. Averaging $M$ independent half-time unitary features first,
then taking their Gram matrix, gives a positive semidefinite estimator.
Within each sampled tuple the same unitary is applied to every parameter
index $i$; different tuples use independent draws.
If $S=\sum_i\|A_i\|_{\rho,1/2}^2$, its operator error is at most
$2\sqrt S\,e_M+e_M^2$ with probability at least $1-\delta$, where
\[
 e_M=\sqrt{S/M}\bigl(1+\sqrt{2\log(1/\delta)}\bigr).
\]
Indeed, the expected squared feature error is at most $S/M$, and bounded
differences apply to its Hilbert norm. The expectation of the Gram estimator
is $(1-1/M)\mathsf G_r+\mathsf G_0/M$, where
$(\mathsf G_0)_{ij}=\langle A_i,A_j\rangle_{\rho,1/2}$. Its bias is explicit. This protocol
assumes access to features or their inner products; it does not replace the
measurement cost of obtaining them on quantum hardware.

\subsection{From Lee--Yang zeros to reusable channel constructions}

The preceding construction depends on a positive zero measure, not on the
Bessel differential equation. Suppose an even real entire function has
the normalized Laguerre--P\'olya factorization
\[
 F(z)=e^{-az^2}\prod_k(1-z^2/r_k^2),\qquad
 a\ge0,\quad r_k>0,\quad\sum_kr_k^{-2}<\infty.
\]
Then
\[
 \frac1{F(ix)}=e^{-ax^2}\prod_k(1+x^2/r_k^2)^{-1}
\]
is an infinitely divisible characteristic function. Substitution of
$\ad_G$ gives a random-unitary semigroup with heat drift $a$ and Gamma
scales $r_k^{-2}$. The same quadrature theorem compiles it. When $F$ is
itself a characteristic function, $(F,1/F(i\,\cdot))$ is a van Dantzig
pair; the relation to Lee--Yang functions is developed in
\citet{KonstantopoulosPatieSarkar2024}.

This supplies a precise route from a zero-location theorem to an operator
algorithm: real zeros determine positive scales, the scales determine a
clock, and quadrature yields a short positive approximation with a local
commutator error. It invokes no unproved assertion about the zeros of a
particular function. The zero-location hypothesis is a named input to the
construction and can be certified independently in finite examples.
